\section*{Week 4}
\subsection*{Exercise 4.1}
Find all the solutions to the system of congruences:
$$ \begin{cases} x \equiv 137 \pmod{423} \\ x \equiv 87 \pmod{191} \end{cases} $$
Use the method by hand (substitution).

\begin{solution}
The system has the form $x \equiv a_1 \pmod{m_1}$ and $x \equiv a_2 \pmod{m_2}$, with $m_1=423$, $m_2=191$, $a_1=137$, and $a_2=87$.

\paragraph{Step 1: Check Solvability and Coprimality}
We first check if $\gcd(423, 191) = 1$ using the Euclidean Algorithm:
\begin{align*}
    423 &= 191 \cdot 2 + 41 \\
    191 &= 41 \cdot 4 + 27 \\
    41 &= 27 \cdot 1 + 14 \\
    27 &= 14 \cdot 1 + 13 \\
    14 &= 13 \cdot 1 + 1 \\
    13 &= 1 \cdot 13 + 0
\end{align*}
Since the last non-zero remainder is $1$, $\gcd(423, 191) = 1$. A unique solution modulo $M = 423 \cdot 191$ exists [Thm 8.1: Chinese Remainder Theorem].

\paragraph{Step 2: Substitution and Reduction}
From the first congruence, $x$ must be of the form:
$$ x = 137 + 423k \quad \text{for some } k \in \mathbb{Z} \quad [\text{Def: Congruence}]$$
Substitute this into the second congruence:
$$ 137 + 423k \equiv 87 \pmod{191} $$
$$ 423k \equiv 87 - 137 \pmod{191} $$
$$ 423k \equiv -50 \pmod{191} $$
To simplify, we reduce the coefficient $423$ modulo $191$ [Prop 6.1: Compatibility, for multiplication]. Using the first step of the Euclidean Algorithm:
$$ 423 = 191 \cdot 2 + 41 \implies 423 \equiv 41 \pmod{191} $$
Also, $-50 \equiv -50 + 191 \equiv 141 \pmod{191}$.
The linear congruence for $k$ becomes:
$$ 41k \equiv 141 \pmod{191} \quad [\text{Reduced congruence for } k] $$

\paragraph{Step 3: Solve for $k$}
We need the inverse of $41 \pmod{191}$. We use the intermediate steps from the Euclidean Algorithm in Step 1 to apply the Extended Euclidean Algorithm:
\begin{align*}
    1 &= 14 - 13 \cdot 1 \\
    1 &= 14 - (27 - 14 \cdot 1) \cdot 1 & [\text{Substitute } 13 = 27 - 14 \cdot 1] \\
    1 &= 14 \cdot 2 - 27 \cdot 1 \\
    1 &= (41 - 27 \cdot 1) \cdot 2 - 27 \cdot 1 & [\text{Substitute } 14 = 41 - 27 \cdot 1] \\
    1 &= 41 \cdot 2 - 27 \cdot 3 \\
    1 &= 41 \cdot 2 - (191 - 41 \cdot 4) \cdot 3 & [\text{Substitute } 27 = 191 - 41 \cdot 4] \\
    1 &= 41 \cdot 2 - 191 \cdot 3 + 41 \cdot 12 \\
    1 &= 41 \cdot 14 - 191 \cdot 3
\end{align*}
Thus, $41(14) \equiv 1 \pmod{191}$. The inverse is $41^{-1} \equiv 14 \pmod{191}$.

Now, multiply the congruence for $k$ by the inverse $14$:
$$ 14 \cdot (41k) \equiv 14 \cdot 141 \pmod{191} $$
$$ k \equiv 1974 \pmod{191} $$
We reduce $1974$ modulo $191$:
$$ 1974 = 191 \cdot 10 + 64 $$
$$ k \equiv 64 \pmod{191} \quad [\text{Solution for } k] $$
This means $k = 64 + 191l$ for some $l \in \mathbb{Z}$.

\paragraph{Step 4: Final Solution for $x$}
Substitute the expression for $k$ back into $x = 137 + 423k$:
\begin{align*}
    x &= 137 + 423(64 + 191l) \\
    x &= 137 + 423 \cdot 64 + 423 \cdot 191l \\
    x &= 137 + 27072 + (423 \cdot 191)l \\
    x &= 27209 + 80793l
\end{align*}
The modulus for the solution is $M = 423 \cdot 191 = 80793$.
We reduce the particular solution $x_0 = 27209$ modulo $80793$. Since $27209 < 80793$, the solution is $x \equiv 27209 \pmod{80793}$.
\textbf{Result:} $x \equiv 27209 \pmod{80793}$.
\end{solution}

---

\subsection*{Exercise 4.2 (Smallest Secret Number)}
An adventurer has the following conditions for a secret number $x$:
\begin{itemize}
\item $x$, when divided by 3, leaves a remainder of 2.
\item $x$, when divided by 5, leaves a remainder of 3.
\item $x$, when divided by 7, leaves a remainder of 4.
\end{itemize}
Find the smallest secret number that satisfies all three conditions.

\begin{solution}
This system of congruences is:
$$ \begin{cases} x \equiv 2 \pmod{3} \quad (a_1=2, m_1=3) \\ x \equiv 3 \pmod{5} \quad (a_2=3, m_2=5) \\ x \equiv 4 \pmod{7} \quad (a_3=4, m_3=7) \end{cases} $$
Since $3, 5, 7$ are pairwise coprime, a unique solution exists modulo $M = 3 \cdot 5 \cdot 7 = 105$ [Thm 8.1: Chinese Remainder Theorem]. We use the constructive method.

\paragraph{Step 1: Compute $M$ and $N_i$}
$$ M = m_1 m_2 m_3 = 3 \cdot 5 \cdot 7 = 105 $$
$$ N_1 = M/m_1 = 105/3 = 35 $$
$$ N_2 = M/m_2 = 105/5 = 21 $$
$$ N_3 = M/m_3 = 105/7 = 15 $$

\paragraph{Step 2: Find the Modular Inverses $y_i$ (Bézout coefficients)}
We solve $N_i y_i \equiv 1 \pmod{m_i}$ for $y_i$.
\begin{itemize}
\item $i=1$: $35 y_1 \equiv 1 \pmod{3}$. Since $35 = 3 \cdot 11 + 2$, $35 \equiv 2 \pmod{3}$.
    $$ 2 y_1 \equiv 1 \pmod{3} $$
    Since $2 \cdot 2 = 4 \equiv 1 \pmod{3}$, the inverse is $y_1 = 2$.
    
    \item $i=2$: $21 y_2 \equiv 1 \pmod{5}$. Since $21 = 5 \cdot 4 + 1$, $21 \equiv 1 \pmod{5}$.
    $$ 1 y_2 \equiv 1 \pmod{5} $$
    The inverse is $y_2 = 1$.
    
    \item $i=3$: $15 y_3 \equiv 1 \pmod{7}$. Since $15 = 7 \cdot 2 + 1$, $15 \equiv 1 \pmod{7}$.
    $$ 1 y_3 \equiv 1 \pmod{7} $$
    The inverse is $y_3 = 1$.
\end{itemize}

\paragraph{Step 3: Construct the Solution $x_0$}
The particular solution is constructed as $x_0 = \sum_{i=1}^3 a_i N_i y_i$ [Thm 8.1 Proof: Existence (Construction)].
$$ x_0 = a_1 N_1 y_1 + a_2 N_2 y_2 + a_3 N_3 y_3 $$
$$ x_0 = (2 \cdot 35 \cdot 2) + (3 \cdot 21 \cdot 1) + (4 \cdot 15 \cdot 1) $$
$$ x_0 = 140 + 63 + 60 = 263 \quad [\text{Intermediate sum}] $$

\paragraph{Step 4: Find the Smallest Positive Solution}
The solution $x$ is unique modulo $M=105$. We reduce $x_0 = 263$ modulo $105$:
$$ 263 = 105 \cdot 2 + 53 $$
$$ x \equiv 53 \pmod{105} \quad [\text{Solution modulo } M] $$
The smallest secret number is the smallest positive residue, which is $53$.
\textbf{Result:} The smallest secret number is $x=53$.
\end{solution}

---

\subsection*{Exercise 4.6}
Find all the solutions to the system of congruences:
$$ \begin{cases} x \equiv 5 \pmod{9} \\ x \equiv 6 \pmod{10} \\ x \equiv 7 \pmod{11} \end{cases} $$
Use the constructive method of the proof of the Theorem (CRT).

\begin{solution}
The system is defined by moduli $m_i = \{9, 10, 11\}$ and residues $a_i = \{5, 6, 7\}$. Since $\gcd(9, 10) = 1$, $\gcd(9, 11) = 1$, and $\gcd(10, 11) = 1$, the moduli are pairwise coprime, and a unique solution exists modulo $M=9 \cdot 10 \cdot 11$ [Thm 8.1: Chinese Remainder Theorem].

\paragraph{Step 1: Compute $M$ and $N_i$}
$$ M = 9 \cdot 10 \cdot 11 = 990 $$
$$ N_1 = M/m_1 = 990/9 = 110 $$
$$ N_2 = M/m_2 = 990/10 = 99 $$
$$ N_3 = M/m_3 = 990/11 = 90 \quad [\text{The } N_i \text{ values are } \{110, 99, 90\}] \text{}$$

\paragraph{Step 2: Find the Modular Inverses $y_i$}
We solve $N_i y_i \equiv 1 \pmod{m_i}$ for $y_i$.
\begin{itemize}
    \item $i=1$: $110 y_1 \equiv 1 \pmod{9}$. Since $110 = 9 \cdot 12 + 2$, $110 \equiv 2 \pmod{9}$.
    $$ 2 y_1 \equiv 1 \pmod{9} $$
    Since $2 \cdot 5 = 10 \equiv 1 \pmod{9}$, the inverse is $y_1 = 5$.
    
    \item $i=2$: $99 y_2 \equiv 1 \pmod{10}$. Since $99 = 10 \cdot 9 + 9$, $99 \equiv 9 \pmod{10}$.
    $$ 9 y_2 \equiv 1 \pmod{10} $$
    Since $9 \equiv -1 \pmod{10}$, and $(-1)(-1) = 1$, the inverse is $y_2 = -1 \equiv 9 \pmod{10}$.
    
    \item $i=3$: $90 y_3 \equiv 1 \pmod{11}$. Since $90 = 11 \cdot 8 + 2$, $90 \equiv 2 \pmod{11}$.
    $$ 2 y_3 \equiv 1 \pmod{11} $$
    Since $2 \cdot 6 = 12 \equiv 1 \pmod{11}$, the inverse is $y_3 = 6$.
\end{itemize}
\paragraph{Step 3: Construct the Solution $x_0$}
$$ x_0 = \sum_{i=1}^3 a_i N_i y_i $$
$$ x_0 = (5 \cdot 110 \cdot 5) + (6 \cdot 99 \cdot 9) + (7 \cdot 90 \cdot 6) $$
$$ x_0 = 2750 + 5346 + 3780 = 11876 \quad [\text{Intermediate sum}] $$

\paragraph{Step 4: Find the Smallest Positive Solution}
The solution $x$ is unique modulo $M=990$. We reduce $x_0 = 11876$ modulo $990$:
$$ 11876 \div 990 \approx 12 \quad [\text{Quotient is } 12] \text{}$$
$$ r = 11876 - 990 \cdot 12 = 11876 - 11880 = -4 $$
$$ x \equiv -4 \equiv -4 + 990 \equiv 986 \pmod{990} \quad [\text{Smallest positive solution}] $$
\textbf{Result:} The set of all solutions is $x \equiv 986 \pmod{990}$.
\end{solution}